\documentclass[conference]{IEEEtran}
\usepackage[utf8]{inputenc}
\usepackage{amsmath}
\usepackage{booktabs}
\usepackage{verbatim}
\usepackage{graphicx}

\title{Labwork 2: Logistic Regression}
\author{\IEEEauthorblockN{Nguyen Phong Chau}
\IEEEauthorblockA{Student ID 2540008\\
University of Science and Technology of Hanoi}
}

\begin{document}

\onecolumn

\maketitle

\begin{abstract}
This lab work reports a logistic regression experiment using gradient descent on a dataset to recommend loan approvals. The study also evaluates the convergence behaviour and final model performance.
\end{abstract}

\section{Implementation}

The experiment uses initial value $w=(0.0, -1.0, 1.0)$, learning rate $\eta=1 \times 10^{-4}$, and convergence threshold $\epsilon=10^{-6}$. The training loop iteratively updates the weights until the change in loss is less than $\epsilon$.

The training algorithm is implemented as follows:

\begin{verbatim}
import math


def predict(x, w) -> float:
    x_with_bias = [1.0] + list(x)
    return sum((w_i * x_i for w_i, x_i in zip(w, x_with_bias)))


def sigmoid(z) -> float:
    return 1 / (1 + math.exp(-z))


def loss(x, y, w) -> float:
    y_pred = [predict(x_i, w) for x_i in x]
    return (
        sum(
            math.log(1 + math.exp(y_pred_i)) - y_i * y_pred_i
            for y_i, y_pred_i in zip(y, y_pred)
        )
        / size
    )


def grad_descent(x, y, w, lr):
    y_pred = [predict(x_i, w) for x_i in x]
    x = [[1.0] + list(x_i) for x_i in x]
    new_w = []
    for j, w_j in enumerate(w):
        w_j = w_j - (
            lr
            * sum(
                x_i[j] * (1 - y_i - sigmoid(-y_pred_i)) for x_i, y_i, y_pred_i in zip(x, y, y_pred)
            )
            / size
        )
        new_w.append(w_j)
    return new_w


def training_loop(weight, lr, epsilon):
    time = 0
    current_loss = loss(x_train, y_train, weight)
    # print("initial loss", current_loss)
    delta = 10

    while delta > epsilon:
        weight = grad_descent(x_train, y_train, weight, lr)
        new_loss = loss(x_train, y_train, weight)
        delta = current_loss - new_loss
        current_loss = new_loss
        # print(time, weight, current_loss, delta)

        time += 1
    return time, weight, current_loss
\end{verbatim}

\section{Results}

The model converged after $81638$ iterations, yielding a final weight of approximately $(-0.2353, 0.9585, -0.1156)$ and a final MSE loss of $0.4979$. 

The final regression line compare to the original data points is shown in Figure~\ref{fig:convergence}.

\begin{figure}[ht]
    \centering
    \includegraphics[width=1\textwidth]{output.png}
    \caption{Compare model prediction with original data points}
    \label{fig:convergence}
\end{figure}

\section{Discussion}

Experimenting different learning rates revealed that $\eta=3 \times 10^{-4}$ provided the fastest convergence time. Convergence time increased faster in both directions when the learning rate was further increased or decreased from this optimal value.

The relation between learning rate and convergence time is shown in the Figure~\ref{fig:lr_convergence}.

\begin{figure}[ht]
    \centering
    \includegraphics[width=1\textwidth]{lr_convergence.png}
    \caption{Convergence time for different learning rates}
    \label{fig:lr_convergence}
\end{figure}


\end{document}